\documentclass[
	11pt
]{beamer}
\usetheme{Madrid}
\useoutertheme{default} % bỏ hiệu ứng shadow
\setbeamertemplate{background canvas}[default] % bỏ bóng đổ hoàn toàn

\setbeamertemplate{footline}{} % bỏ thanh cuối
\setbeamertemplate{navigation symbols}{} % bỏ icon điều hướng


\usepackage[utf8]{inputenc}   % Cho phép nhập Unicode
\usepackage[T1]{fontenc}      % Font encoding hỗ trợ ký tự toán học
\usepackage{lmodern}          % Font hiện đại
\usepackage{textcomp}         % Bổ sung ký tự đặc biệt
\usepackage[vietnamese]{babel} % Nếu dùng tiếng Việt
\graphicspath{{Images/}{./}} % Specifies where to look for included images (trailing slash required)

\usepackage{booktabs} % Allows the use of \toprule, \midrule and \bottomrule for better rules in tables
\usepackage{amsmath} % Để sử dụng các ký hiệu toán học như P[...]
\usepackage{tikz} % Để vẽ các dấu ngoặc và chú thích
\usepackage{svg}
\usetikzlibrary{positioning, shapes.geometric, arrows}
\usepackage{tikz}
\usepackage{pgfplots}
\pgfplotsset{compat=1.18}
\usepackage{pgf-pie}

%----------------------------------------------------------------------------------------
%	SELECT LAYOUT THEME
%----------------------------------------------------------------------------------------



%\usetheme{default}
%\usetheme{AnnArbor}
%\usetheme{Antibes}
%\usetheme{Bergen}
%\usetheme{Berkeley}
%\usetheme{Berlin}
%\usetheme{Boadilla}
%\usetheme{CambridgeUS}
%\usetheme{Copenhagen}
%\usetheme{Darmstadt}
%\usetheme{Dresden}
%\usetheme{Frankfurt}
%\usetheme{Goettingen}
%\usetheme{Hannover}
%\usetheme{Ilmenau}
%\usetheme{JuanLesPins}
%\usetheme{Luebeck}


%\usetheme{Malmoe}
%\usetheme{Marburg}
%\usetheme{Montpellier}
%\usetheme{PaloAlto}
%\usetheme{Pittsburgh}
%\usetheme{Rochester}
%\usetheme{Singapore}
%\usetheme{Szeged}
%\usetheme{Warsaw}




\usefonttheme{default} 
\usepackage{palatino} % Use the Palatino font for serif text


\usepackage[default]{opensans} 

%----------------------------------------------------------------------------------------
%	PRESENTATION INFORMATION
%----------------------------------------------------------------------------------------
\title[BÁO CÁO KẾT QUẢ BÀI TẬP]{BÁO CÁO KẾT QUẢ BÀI TẬP}
\subtitle{Phân loại người sống sót trên tàu Titanic}

\author[Nhóm 4]{%
	\textbf{Giảng viên hướng dẫn:} TS. Đỗ Như Tài \\[0.8em]
	\textbf{Sinh viên thực hiện:}\\[0.3em]
	Trần Đình Hải \\ 
	Huỳnh Nhật Minh \\ 
	Lê Thành Nhân \\ 
	Trần Công Minh
}

\institute[SGU]{Trường Đại học Sài Gòn}
\date{\today}




\begin{document}
	

%----------------------------------------------------------------------------------------
%	TITLE SLIDE
%----------------------------------------------------------------------------------------

\begin{frame}
	\titlepage % Output the title slide, automatically created using the text entered in the PRESENTATION INFORMATION block above
\end{frame}

%----------------------------------------------------------------------------------------
%	TABLE OF CONTENTS SLIDE
%----------------------------------------------------------------------------------------


\begin{frame}
	\frametitle{Presentation Overview} 
	
	\tableofcontents 
\end{frame}

%----------------------------------------------------------------------------------------
%	PRESENTATION BODY SLIDES
%----------------------------------------------------------------------------------------
%==================================================
%================ SLIDE 1 ==================
\section{Định nghĩa vấn đề}
\begin{frame}{Định nghĩa vấn đề}
\begin{itemize}
    \item Mục tiêu: Dự đoán hành khách \textbf{sống sót hay không} trên tàu Titanic.
    \item \textbf{train.csv}: chứa đặc điểm hành khách và nhãn \texttt{Survived (0 = Không, 1 = Có)}.
    \item \textbf{test.csv}: chứa các đặc điểm tương tự, \textbf{không có cột Survived}.
    \item Dữ liệu dùng để huấn luyện và kiểm tra mô hình học máy.
\end{itemize}
\end{frame}

%================ SLIDE 2 ==================
\begin{frame}{Đặc trưng đầu vào và đầu ra}

\begin{itemize}
    \item \textbf{Pclass}: Hạng vé (1–3)
    \item \textbf{Name}: Tên hành khách
    \item \textbf{Sex}: Giới tính (nam/nữ)
    \item \textbf{Age}: Tuổi (năm)
    \item \textbf{SibSp}: Số anh chị em / vợ chồng đi cùng
    \item \textbf{Parch}: Số cha mẹ / con cái đi cùng
    \item \textbf{Ticket}: Số vé
    \item \textbf{Fare}: Giá vé (bảng Anh)
    \item \textbf{Cabin}: Mã phòng hoặc khoang tàu
    \item \textbf{Embarked}: Cảng lên tàu (C/Q/S)
\end{itemize}


%------ Cột phải: sơ đồ minh họa ------

\centering
\begin{tikzpicture}[node distance=1.2cm,>=stealth,thick]
\tikzstyle{box}=[rectangle, rounded corners, draw=blue!50, fill=blue!10,
                 text width=3.2cm, text centered, minimum height=1cm]
\tikzstyle{arrow}=[->, line width=1pt]

\node[box] (input) {Đầu vào\\(Pclass, Sex, Age, Fare, ...)};
\node[box, right=of input, text width=2.8cm] (model) {Mô hình học máy};
\node[box, right=of model, fill=green!10, draw=green!50,text width=2.8cm] (output) {Đầu ra\\Survived (0/1)};

\draw[arrow] (input) -- (model);
\draw[arrow] (model) -- (output);
\end{tikzpicture}

\end{frame}

\section{Phân tích dữ liệu khám phá - EDA}
%================ SLIDE 1 ==================
\begin{frame}{Tổng quan và tính toàn vẹn dữ liệu}
\begin{itemize}
    \item Tập \textbf{train.csv}: 891 dòng, 12 cột  
    \item Biến gồm dạng số và dạng chuỗi  
    \item Thiếu dữ liệu ở: \texttt{Age}, \texttt{Cabin}, \texttt{Embarked}  
    \item \textbf{Age:} 177 thiếu (~20\%), \textbf{Cabin:} 687 thiếu (~77\%), \textbf{Embarked:} 2 thiếu (<1\%)  
    \item Cột mục tiêu: \textbf{Survived (0 = Không, 1 = Có)}  
    \item Không có bản ghi trùng → Dữ liệu duy nhất (unique)  
\end{itemize}
\end{frame}

%================ SLIDE 2 ==================
\begin{frame}{Thống kê mô tả và Phân bố biến phân lớp (Survived)}
	
	\textbf{Thống kê mô tả các biến số:}
	\begin{itemize}
		\item \textbf{Pclass:} Trung bình 2.31 $\rightarrow$ đa số hạng 2–3  
		\item \textbf{Age:} TB $\approx$ 29.7, lệch chuẩn $\approx$ 14.5 (0.42–80)  
		\item \textbf{SibSp/Parch:} Đa số đi một mình, tối đa 8 anh chị em hoặc 6 cha mẹ/con  
		\item \textbf{Fare:} TB $\approx$ 32.2, lệch chuẩn $\approx$ 49.7, có outlier đến 512.33  
	\end{itemize}
	
	\vspace{0.6em}
	
	\textbf{Phân bố biến phân lớp \texttt{Survived}:}
	\begin{itemize}
		\item Dữ liệu hơi \textbf{mất cân bằng (imbalanced)}.
	\end{itemize}
	
	\vspace{0.6em}
	
	\begin{center}
		\begin{tikzpicture}
			\pie[
			text=legend,
			radius=1.8,
			color={red!60, green!60},
			sum=auto
			]{
				62/Không sống sót,
				38/Sống sót
			}
		\end{tikzpicture}
		
		\vspace{0.3em}
		\footnotesize
		\textit{Biểu đồ 1: Phân bố tỉ lệ sống sót trong tập train.}
	\end{center}
	
\end{frame}



\begin{frame}{ Đơn biến (Univariate Plots)}

\begin{figure}
    \centering
    \includegraphics[width=0.55\linewidth]{countplot.png}
    
    
\end{figure}
\begin{itemize}
    \item Số lượng hành khách \textbf{không sống sót (549)} cao hơn đáng kể so với \textbf{sống sót (342)}.
    \item Phần lớn hành khách đã không qua khỏi, chiếm khoảng \textbf{61.6\%} tổng khách.
\end{itemize}

\vspace{1.5em}



\end{frame}

\begin{frame}{ Box and Whisker Plots}


\begin{itemize}
    \item \textbf{Age:} 20–38 tuổi, trung vị 28–30, có outliers lớn tuổi.  
    \item \textbf{SibSp:} 75\% từ 0–1, trung vị 0, vài giá trị ngoại lệ.  
    \item \textbf{Parch:} Đa số 0, có outliers.  
    \item \textbf{Fare:} Lệch phải mạnh, trung vị < 20, outliers > 500.  
    \item \textbf{Fare} có thang đo lớn → các biến khác bị nén ở đáy.  
\end{itemize}


\begin{columns}[c]
    \column{0.5\textwidth}
    \begin{center}
        \includegraphics[width=\linewidth]{output.png}
    \end{center}

    \column{0.5\textwidth}
    \begin{center}
        \includegraphics[width=\linewidth]{output1.png}
    \end{center}
\end{columns}
\end{frame}
\begin{frame}{Phân bố tuổi theo tình trạng sống sót}
	
	\begin{figure}[h]
		\centering
		\includegraphics[width=0.6\linewidth]{age_distribution_survival.png}
		
	\end{figure}
	
	\begin{itemize}
		\item Trẻ em được ưu tiên rõ rệt  
		\item Đỉnh sống sót cao ở nhóm <10  
		\item Người lớn (20–40) tử vong nhiều  
		\item Đây là nhóm tuổi đông nhất  
		\item “Ưu tiên trẻ em” thể hiện rõ  
		\item Nam 20–40 có rủi ro cao nhất  
	\end{itemize}
	
\end{frame}

\begin{frame}{Histogram (Sex vs Survived)}
	
	
	\begin{itemize}
		\item \textbf{Giới tính là yếu tố dự đoán mạnh nhất.}
		
		\item \textbf{Nữ (female):}
		\begin{itemize}
			\item Tỷ lệ sống sót vượt trội.
			\item Số sống sót cao hơn tử vong.
			\item 233/342 người sống sót là nữ (≈68\%).
		\end{itemize}
		
		\item \textbf{Nam (male):}
		\begin{itemize}
			\item Tỷ lệ tử vong rất cao.
			\item 468/549 người tử vong là nam (≈85\%).
		\end{itemize}
		
	\end{itemize}
	
	
	\begin{center}
		\includegraphics[width=0.5\linewidth]{barplot.png}
	\end{center}
	
\end{frame}

\begin{frame}{Phân tích theo hạng vé (Pclass)}
	\small
	\begin{itemize}
		\item Hạng vé càng cao → cơ hội sống càng lớn.  
		\item Hạng 1: sống sót nhiều hơn tử vong.  
		\item Hạng 2: gần cân bằng.  
		\item Hạng 3: tử vong áp đảo.  
		\item Pclass thể hiện rõ ảnh hưởng của địa vị xã hội.
	\end{itemize}
	
	\vspace{0.5em}
	\centering
	\includegraphics[width=0.5\textwidth]{pclass_survival.png}
\end{frame}
\begin{frame}{Phân tích đa biến (Multivariate Analysis)}
	\small
	\begin{itemize}
		\item Hạng vé cao $\Rightarrow$ cơ hội sống cao hơn.  
		\item Người sống sót thường trẻ và trả vé cao hơn.  
		\item Gia đình nhỏ (1–2 người đi cùng) có tỉ lệ sống tốt hơn.  
		\item \textbf{Fare} và \textbf{Pclass} tương quan nghịch rõ rệt.  
		\item \textbf{Age}, \textbf{SibSp}, \textbf{Parch} ít tương quan mạnh.
	\end{itemize}
	
	\vspace{0.5em}
	\centering
	\includegraphics[width=0.5\textwidth]{multivariate_plot.png}
\end{frame}
\begin{frame}{Ma trận tương quan}
	\small
	\begin{itemize}
		\item Không biến nào |r| $>$ 0.8  
		\item Pclass-Fare: tương quan mạnh  
		\item SibSp-Parch: liên hệ nhẹ  
		\item Age: gần như độc lập  
	\end{itemize}
	
	\centering
	\includegraphics[width=0.58\textwidth]{correlation_matrix.png}
\end{frame}


\begin{frame}{Phân bố sống sót}
	\begin{itemize}
		\item Age: Trẻ có lợi thế sống sót  
		\item Fare: Vé cao hạng sang sống sót nhiều  
		\item Pclass: Hạng 1–2 sống sót, hạng 3 tử vong đa số  
		\item SibSp/Parch: có 1–2 người đi cùng sống cao hơn
	\end{itemize}
	
	\centering
	\includegraphics[width=0.5\textwidth]{comparison_boxplots.png}
\end{frame}

\section{Chuẩn bị dữ liệu và xử lý đặc trưng}
\begin{frame}{Tạo đặc trưng mới (Feature Engineering)}
	\small
	\textbf{(1) FamilySize = SibSp + Parch + 1}
	\begin{itemize}
		\item Người đi một mình sống sót rất ít.  
		\item Gia đình nhỏ (2–4 người) có tỉ lệ sống cao nhất.  
		\item Gia đình lớn (>5) khó sống sót do khó hỗ trợ nhau.  
	\end{itemize}
	
	\vspace{0.4em}
	\textbf{(2) Title (Mr, Mrs, Miss, Master, Other)}
	\begin{itemize}
		\item \textbf{Mr:} đông nhất, tỉ lệ tử vong cao.  
		\item \textbf{Mrs, Miss, Master:} sống sót nhiều hơn rõ rệt.  
		\item \textbf{Other:} hiếm, đa phần tử vong.  
	\end{itemize}
\end{frame}

\begin{frame}{Xử lý giá trị thiếu và Mã hoá dữ liệu}
	\small
	\begin{itemize}
		\item \textbf{Age:} điền bằng giá trị trung bình.  
		\item \textbf{Embarked:} thay bằng giá trị xuất hiện nhiều nhất.  
		\item \textbf{Cabin:} xoá do thiếu quá nhiều.  
		\item \textbf{Mã hoá danh mục:} dùng LabelEncoder cho \textit{Sex, Embarked, Title, Family\_cate}.  
		\item Sau mã hoá:  
		\textit{Sex → [female, male]}, \textit{Embarked → [C, Q, S]},  
		\textit{Title → [Master, Miss, Mr, Mrs, Other]},  
		\textit{Family\_cate → [Alone, Small, Medium, Large]}.  
	\end{itemize}
\end{frame}

\begin{frame}{Chuẩn hóa dữ liệu}
	
	\textbf{Phương pháp:}  
	Dữ liệu được chuẩn hóa bằng \textbf{StandardScaler} và \textbf{MinMaxScaler}  
	để đưa các đặc trưng về cùng thang đo, tránh sai lệch do đơn vị khác nhau.
	
	\vspace{0.8em}
	\centering
	\begin{tabular}{|l|c|c|c|}
		\hline
		\textbf{Biến} & \textbf{Phạm vi gốc} & \textbf{MinMaxScaler} & \textbf{StandardScaler} \\
		\hline
		Age & 0.42 -- 80.00 & [0, 1] & [-2.2, 2.5] \\
		Fare & 0.00 -- 512.33 & [0, 1] & [-0.6, 9.8] \\
		Parch & 0 -- 6 & [0, 1] & [-0.5, 6.0] \\
		SibSp & 0 -- 8 & [0, 1] & [-0.5, 7.0] \\
		FamilySize & 1 -- 11 & [0, 1] & [-0.5, 6.0] \\
		\hline
	\end{tabular}
	
	\vspace{0.6em}
	\small{
		Các giá trị min và max được tính tự động từ tập huấn luyện.  
		\textbf{StandardScaler} đưa dữ liệu về trung bình 0 và độ lệch chuẩn 1,  
		nhưng do tồn tại nhiều outlier (đặc biệt ở \textit{Fare}),  
		một số giá trị vượt quá khoảng [-3, 3].
	}
	
\end{frame}


\section{Huấn luyện mô hình}




\begin{frame}{Đánh Giá -- StandardScaler}
	\begin{columns}[T,onlytextwidth]
		\column{0.45\textwidth}
		\textbf{Dữ liệu:} \texttt{df\_standard.xlsx}
		\vspace{0.4em}
		\textbf{Chiến lược:}
		K-Fold CV ($k=5$, Stratified)
		$\Rightarrow$ Metric: \textbf{Accuracy (Mean $\pm$ Std)}
		\vspace{0.8em}
		\textbf{Kích thước:}
		$X_{train}$: (891, 10 features) \\
		$y_{train}$: (891,) \\
		Mỗi fold: $\sim$712 train, $\sim$179 valid
		
		\column{0.55\textwidth}
		\centering
		\begin{tikzpicture}[
			node distance=0.85cm and 0.6cm,
			every node/.style={font=\scriptsize, align=center},
			stage/.style={
				rectangle,
				rounded corners=4pt,
				draw=blue!70,
				fill=blue!5,
				thick,
				minimum width=3.3cm,
				minimum height=0.6cm
			},
			arrow/.style={->, thick, >=stealth, blue!70}
			]
			\node[stage] (s1) {1. Chia dữ liệu \\ $X_{train}, y_{train}$ \\ thành 5 fold};
			\node[stage, below=of s1] (s2) {2. Mỗi fold: \\ Train 4 phần \\ Valid 1 phần};
			\node[stage, below=of s2] (s3) {3. Tính Accuracy \\ (Mean $\pm$ Std)};
			\node[stage, below=of s3] (s4) {4. Huấn luyện lại \\ toàn bộ dữ liệu};
			\draw[arrow] (s1) -- (s2);
			\draw[arrow] (s2) -- (s3);
			\draw[arrow] (s3) -- (s4);
			\node[
			draw=blue!50,
			fill=blue!10,
			rounded corners=3pt,
			inner sep=3pt,
			text width=4.4cm,
			below=0.4cm of s4,
			font=\tiny
			]
			{Quy trình đánh giá với \textbf{StandardScaler} \\ chuẩn hóa theo z-score};
		\end{tikzpicture}
	\end{columns}
\end{frame}

\begin{frame}{Kết Quả Baseline -- StandardScaler}
	\centering
	
	\begin{tabular}{lccc}
		\toprule
		\textbf{Mô hình} & \textbf{Acc} & \textbf{Std} & \textbf{Rank} \\
		\midrule
		SVM     & \textbf{0.8238} & $\pm$0.0175 & 1 \\
		NB      & 0.8014 & $\pm$0.0213 & 2 \\
		KNN     & 0.8013 & $\pm$0.0171 & 3 \\
		LDA     & 0.8002 & $\pm$0.0327 & 4 \\
		LogReg  & 0.7890 & $\pm$0.0204 & 5 \\
		CART    & 0.7744 & $\pm$0.0188 & 6 \\
		\bottomrule
	\end{tabular}
	
	\vspace{0.4em}
	
	\begin{itemize}
		\item SVM cao nhất (0.8238), ổn định, overfit +2.84\%
		\item NB, KNN $\approx$80.1\%; LDA, LogReg ổn
		\item CART thấp nhất (0.7744), overfit +4.16\%
		\item SVM vượt 1–5\%
	\end{itemize}
\end{frame}

\begin{frame}{Phân Tích Boxplot -- StandardScaler}
	\centering
	\includegraphics[width=0.9\textwidth]{boxplot_baseline_standardscaler.png}
	
	\vspace{0.4em}
	
	\begin{itemize}
		\item SVM: Hiệu suất cao, ổn định tốt
		\item KNN, LDA: Hiệu suất khá, ổn định
		\item NB: Ổn định, accuracy trung bình
		\item LR: Accuracy thấp hơn, ổn định
		\item CART: Thấp nhất, biến thiên cao
		\item Kết luận: SVM vượt trội nhất
	\end{itemize}
\end{frame}



\begin{frame}{Tuning \& Kết Quả Final -- StandardScaler}
	\centering
	
	\begin{tabular}{lcc}
		\toprule
		\textbf{Mô hình} & \textbf{CV Acc} & \textbf{Train Acc} \\
		\midrule
		LogReg  & 0.8070 & 0.8070 \\
		SVM     & \textbf{0.8283} & 0.8339 \\
		\bottomrule
	\end{tabular}
	
	\vspace{0.4em}
	
	\begin{itemize}
		\item LogReg: \texttt{C=1.0, l1, balanced}
		\item SVM: \texttt{C=200, gamma=0.01}
		\item SVM $\uparrow$ +0.45pp vs baseline
		\item Chênh Train–CV: +0.56\% → ổn định
		\item \textbf{Chọn SVM, LR làm final}
	\end{itemize}
\end{frame}




\begin{frame}{\textbf{Kết Quả Test – StandardScaler}}
	
	\centering
	
	\begin{tabular}{lcc}
		\toprule
		\textbf{Mô hình} & \textbf{Not Survived} & \textbf{Survived} \\
		\midrule
		\textbf{LR} & 127 & 291 \\
		\textbf{SVM} & 222 & 196 \\
		\bottomrule
	\end{tabular}
	
	\vspace{0.6em}
	
	\begin{itemize}
		\item \textbf{SVM} dự đoán cân bằng hơn.  
		\item Giảm lệch lớp, ổn định hơn LR.  
		\item \textbf{LR} thiên về lớp sống sót.  
		\item Chốt: \textbf{SVM + StandardScaler} hiệu quả nhất.
	\end{itemize}
	
\end{frame}






\begin{frame}{Đánh Giá -- MinMaxScaler}
	\begin{columns}[T,onlytextwidth]
		\column{0.45\textwidth}
		\textbf{Dữ liệu:} \texttt{df\_minmax.xlsx}
		\vspace{0.4em}
		\textbf{Chiến lược:}
		K-Fold CV ($k=5$, Stratified)
		$\Rightarrow$ Metric: \textbf{Accuracy (Mean $\pm$ Std)}
		\vspace{0.8em}
		\textbf{Kích thước:}
		$X_{train}$: (891, 10 features) \\
		$y_{train}$: (891,) \\
		Mỗi fold: $\sim$712 train, $\sim$179 valid
		
		\column{0.55\textwidth}
		\centering
		\begin{tikzpicture}[
			node distance=0.85cm and 0.6cm,
			every node/.style={font=\scriptsize, align=center},
			stage/.style={
				rectangle,
				rounded corners=4pt,
				draw=teal!70,
				fill=teal!8,
				thick,
				minimum width=3.3cm,
				minimum height=0.6cm
			},
			arrow/.style={->, thick, >=stealth, teal!70}
			]
			\node[stage] (s1) {1. Chia dữ liệu \\ $X_{train}, y_{train}$ \\ thành 5 fold};
			\node[stage, below=of s1] (s2) {2. Mỗi fold: \\ Train 4 phần \\ Valid 1 phần};
			\node[stage, below=of s2] (s3) {3. Tính Accuracy \\ (Mean $\pm$ Std)};
			\node[stage, below=of s3] (s4) {4. Huấn luyện lại \\ toàn bộ dữ liệu};
			\draw[arrow] (s1) -- (s2);
			\draw[arrow] (s2) -- (s3);
			\draw[arrow] (s3) -- (s4);
			\node[
			draw=teal!50,
			fill=teal!10,
			rounded corners=3pt,
			inner sep=3pt,
			text width=4.4cm,
			below=0.4cm of s4,
			font=\tiny
			]
			{Quy trình đánh giá với \textbf{MinMaxScaler} \\ chuẩn hóa dữ liệu về khoảng $[0,1]$};
		\end{tikzpicture}
	\end{columns}
\end{frame}








\begin{frame}{Kết Quả Baseline -- MinMaxScaler}
	\centering
	
	\begin{tabular}{lccc}
		\toprule
		\textbf{Mô hình} & \textbf{Acc} & \textbf{Std} & \textbf{Rank} \\
		\midrule
		SVM     & \textbf{0.8317} & $\pm$0.0117 & 1 \\
		LR      & 0.8047 & $\pm$0.0215 & 2 \\
		KNN     & 0.8036 & $\pm$0.0218 & 3 \\
		LDA     & 0.7991 & $\pm$0.0224 & 4 \\
		NB      & 0.7969 & $\pm$0.0080 & 5 \\
		CART    & 0.7755 & $\pm$0.0291 & 6 \\
		\bottomrule
	\end{tabular}
	
	\vspace{0.4em}
	
	\begin{itemize}
		\item SVM cao nhất (0.8317), ổn định nhất (±0.0117)
		\item LR, KNN $\approx$80.4\%; LDA, NB ổn
		\item CART thấp nhất (0.7755), biến động cao
		\item MinMax > Standard: SVM +0.79\%, CART +0.14\%
	\end{itemize}
\end{frame}


\begin{frame}{Boxplot Baseline -- MinMaxScaler}
	\centering
	\includegraphics[width=0.9\textwidth]{boxplot_baseline_minmaxscaler.png}
	
	\vspace{0.4em}
	
	\begin{itemize}
		\item SVM: Median 0.8258, hộp hẹp, ổn định
		\item LR: Median 0.8034, ổn định tốt
		\item KNN, LDA: $\approx$0.79, ổn định TB
		\item NB: Ổn định, acc trung bình
		\item CART: Median 0.7584, biến thiên cao
		\item MinMax > Standard: SVM +0.58\%, LR +2.54\%
	\end{itemize}
\end{frame}





\begin{frame}{Tuning \& Kết Quả Final -- MinMaxScaler}
	\centering
	
	\begin{tabular}{lccc}
		\toprule
		\textbf{Mô hình} & \textbf{CV Acc} & \textbf{Train Acc} & \textbf{Gap} \\
		\midrule
		KNN     & 0.8114 & 0.8496 & 3.82\% \\
		SVM     & \textbf{0.8249} & 0.8305 & \textbf{0.56\%} \\
		\bottomrule
	\end{tabular}
	
	\vspace{0.4em}
	
	\begin{itemize}
		\item KNN: \texttt{n=7}
		\item SVM: \texttt{C=1, gamma=auto}
		\item SVM > KNN +1.35\%, ổn định hơn
		\item Gap nhỏ → ít overfit
		\item \textbf{Chọn SVM, KNN làm final}
		\item StandardScaler > MinMax
	\end{itemize}
\end{frame}

\begin{frame}{\textbf{Kết Quả Test – MinMaxScaler}}
	
	\centering
	\small
	\begin{tabular}{lcc}
		\toprule
		\textbf{Mô hình} & \textbf{Not Survived} & \textbf{Survived} \\
		\midrule
		\textbf{SVM} & 225 & 193 \\
		\textbf{KNN (Backup)} & 218 & 200 \\
		\bottomrule
	\end{tabular}
	
	\vspace{0.6em}
	
	\begin{itemize}
		\item \textbf{SVM} dự đoán cân bằng, ổn định.  
		\item \textbf{KNN} thiên nhẹ về lớp sống sót.  
		\item Cả hai có phân bố gần tương đương.  
		\item Chốt: \textbf{SVM + MinMaxScaler} hiệu quả hơn.
	\end{itemize}
	
\end{frame}


\begin{frame}{So Sánh 2 Phương Pháp Chuẩn Hóa}
	
	\centering
	\small
	\begin{tabular}{lcc}
		\toprule
		\textbf{Tiêu chí} & \textbf{StandardScaler (Model 1)} & \textbf{MinMaxScaler (Model 2)} \\
		\midrule
		Data & df\_standard.xlsx & df\_minmax.xlsx \\
		Baseline Best & SVM: 0.8238 $\pm$ 0.0175 & \textbf{SVM: 0.8317 $\pm$ 0.0117} \\
		Baseline Median & SVM: 0.8156 & SVM: 0.8258 \\
		Mô hình Tuning & LR + SVM & KNN + SVM \\
		Tuning Best (CV) & \textbf{SVM: 0.8283} & SVM: 0.8249 \\
		Tuning Train Acc & 0.8339 & 0.8305 \\
		Gap (Overfitting) & 0.56\% & 0.56\% \\
		Kết luận &  Ổn định tốt &  Ổn định tốt \\
		\bottomrule 
	\end{tabular}
	
	\vspace{0.6em}
	
	\begin{itemize}
		\item \textbf{MinMaxScaler:} Hiệu suất baseline cao hơn, ổn định khi chưa tuning.
		\item \textbf{StandardScaler:} Cải thiện hiệu quả hơn sau khi tuning (SVM đạt 82.83\%).
		\item Cả hai phương pháp đều ổn định và cho kết quả gần tương đương.
	\end{itemize}
	
\end{frame}

	
\begin{frame}{Kết Luận }
	\begin{itemize}
		\item Model 1: SVM + StandardScaler \quad \alert{Kaggle: 0.72248}
		\item Model 2: SVM + MinMaxScaler   \quad \alert{Kaggle: 0.71770}
		\item Cả hai ổn định, gap < 1\%
		\item Nộp chính: Model 1 (GridSearchCV)
		\item Nộp phụ:   Model 2 (RandomizedSearchCV)
		\item Hướng tới: XGBoost + Cabin + Age GAN
	\end{itemize}
\end{frame}










\begin{frame}
\centering
\Huge \textbf{Cảm ơn thầy và các bạn!}
\end{frame}

\end{document}
